% =============================================================================
% SECTION II — Benchmark Factor Models
% =============================================================================


% ── SEMAPHORE ──────────────────────────────────────────────────────────────────

This section tests whether the alpha of the three strategies survives
the canonical benchmark spanning regressions of the fixed-income
arbitrage literature---\citet{duarte2007risk}, \citet{fung2004hedge},
and \citet{brooks2020active}---and whether the residual alpha varies
with intermediary funding stress, as the slow-moving capital framework
of \citet{duffie2010presidential} predicts.
Section~\ref{sec:bench_specs} describes the three specifications and
the construction of their euro-denominated counterparts.
Section~\ref{sec:bench_results} presents the regression results and
three complementary evaluations: the alpha by stress regime, the
conditional alpha model of \citet{christopherson1998conditioning}, and
an out-of-sample re-estimation in which the benchmark hedge is formed
from past-only loadings.


% ── II.A BENCHMARK SPECIFICATIONS ─────────────────────────────────────────────

\subsection{Benchmark Specifications}
\label{sec:bench_specs}

Three benchmark specifications are widely used in the empirical
literature on fixed-income arbitrage and active fixed-income
management, and they provide the natural reference points against
which our strategies should be evaluated.

\citet{duarte2007risk}, in the paper that motivates our research
design, test whether fixed-income arbitrage excess returns are
compensation for exposures to equity, banking-sector, and
term-structure risk:
\begin{equation}\label{eq:duarte}
\begin{split}
  r_{i,t} \;=\; \alpha_i
  &+ \beta_1\,\textit{Mkt}_t + \beta_2\,\textit{SMB}_t
   + \beta_3\,\textit{HML}_t + \beta_4\,\textit{UMD}_t + \beta_5\,\textit{RS}_t + \beta_6\,\textit{RI}_t\\
  &+ \beta_7\,\textit{RB}_t + \beta_8\,\textit{R2}_t + \beta_9\,\textit{R5}_t
   + \beta_{10}\,\textit{R10}_t + \varepsilon_{i,t},
\end{split}
\end{equation}
where $\textit{Mkt}$, $\textit{SMB}$, $\textit{HML}$, and
$\textit{UMD}$ are the Fama--French--Carhart equity factors (market,
size, value, and momentum), $\textit{RS}$ is the S\&P bank stock index
excess return, $\textit{RI}$ and $\textit{RB}$ are A/BAA-rated
industrial and bank bond index excess returns, and $\textit{R2}$,
$\textit{R5}$, $\textit{R10}$ are the CRSP Fama bond portfolios at
two-, five-, and ten-year maturities.

\citet{fung2004hedge} propose a seven-factor \emph{asset-based style}
model that tests whether hedge fund returns are explained by equity,
fixed-income, and trend-following exposures:
\begin{equation}\label{eq:funghsieh}
\begin{split}
  r_{i,t} \;=\; \alpha_i
  &+ \beta_1\,\textit{S\&P}_t + \beta_2\,\textit{SC\text{-}LC}_t
   + \beta_3\,\textit{BdOpt}_t + \beta_4\,\textit{FXOpt}_t \\
  &+ \beta_5\,\textit{ComOpt}_t
   + \beta_6\,\textit{10Y}_t
   + \beta_7\,\textit{CredSpr}_t + \varepsilon_{i,t},
\end{split}
\end{equation}
where $\textit{S\&P}$ is the S\&P~500 excess return,
$\textit{SC\text{-}LC}$ is the small-minus-large-cap total return,
$\textit{BdOpt}$, $\textit{FXOpt}$, and $\textit{ComOpt}$ are the
three trend-following factors constructed as lookback straddle returns
on bond, currency, and commodity futures, $\textit{10Y}$ is the excess
return on a 10-year government bond total-return index, and
$\textit{CredSpr}$ is the return spread between BAA-rated corporate
bonds and 10-year government bonds.

\citet{brooks2020active} take a different angle, decomposing the
returns of active fixed-income managers into eight \emph{traditional
risk premia} organised along three channels of risk-taking---duration,
credit, and volatility selling:
\begin{equation}\label{eq:brooks}
\begin{split}
  r_{i,t} \;=\; \alpha_i
  &+ \beta_1\,\textit{Term}_t + \beta_2\,\textit{GlTerm}_t
   + \beta_3\,\textit{GlAgg}_t + \beta_4\,\textit{InflLink}_t 
   + \beta_5\,\textit{CorpCred}_t \\
  &+ \beta_6\,\textit{EmDebt}_t + \beta_7\,\textit{EmCurr}_t 
  + \beta_8\,\textit{USTVol}_t + \varepsilon_{i,t},
\end{split}
\end{equation}
where $\textit{Term}$, $\textit{GlTerm}$, $\textit{GlAgg}$, and
$\textit{InflLink}$ are the total returns on the Treasury, the
currency-hedged global Treasury, global aggregate, and global
inflation-linked bond indices, each in excess of cash,
$\textit{CorpCred}$ is a high-yield corporate bond index in excess of
cash, $\textit{EmDebt}$ is the excess return on a hard-currency
emerging-market debt index, $\textit{EmCurr}$ is the return on an
equal-weighted basket of emerging-market currencies, and
$\textit{USTVol}$ is the return on delta-hedged straddles on 10-year
Treasury futures.

Two ingredients of the original specifications are replaced with
tradable counterparts. In \citet{fung2004hedge}, the bond and credit
factors are yield and spread \emph{changes}; a yield change is not the
return on an investable position, so we use the corresponding
total-return series (the $\textit{10Y}$ and $\textit{CredSpr}$
definitions above). In \citet{brooks2020active},
$\textit{CorpCred}$ blends high yield with leveraged loans; loan
indices are quote-based and not investable at the monthly horizon, so
we retain the high-yield leg only. With every regressor a realised
excess return, the intercept in each specification keeps its
interpretation as a Jensen alpha.

Each benchmark is estimated in two parallel panels. The first uses the
original U.S.-denominated factors as specified by the respective
authors, which is the natural way to replicate each framework as it
was first published. The second uses euro equivalents, since all three
strategies trade European instruments and may therefore be better
explained by European factor exposures. The euro panel is the primary
specification and is reported in the main text; the US-factor
estimates are reported in Appendix~\ref{app:benchmark_us}.

For the euro panel, the Fama--French factors---used directly in
\citet{duarte2007risk} and indirectly in \citet{fung2004hedge}---are
sourced from the Kenneth French data library; because these are
reported in USD, we convert them to EUR following
\citet{gluck2021currency}. For long factors (market excess return):
\begin{equation}\label{eq:gluck_long}
  F_{t}^{\textsc{eur}}
  \;=\; \frac{1 + F_{t}^{\textsc{usd}} + r_{f,t}^{\textsc{usd}}}
             {1 + r_{t}^{\textsc{usd/eur}}}
        \;-\; 1 \;-\; r_{f,t}^{\textsc{eur}},
\end{equation}
and for long-short factors (SMB, HML, UMD):
\begin{equation}\label{eq:gluck_ls}
  \mathit{LS}_{t}^{\textsc{eur}}
  \;=\; \frac{\mathit{LS}_{t}^{\textsc{usd}}}
             {1 + r_{t}^{\textsc{usd/eur}}},
\end{equation}
where $F_{t}^{\textsc{usd}}$ and $\mathit{LS}_{t}^{\textsc{usd}}$ are
the original factors in US dollars, $r_{f,t}^{\textsc{usd}}$ and
$r_{f,t}^{\textsc{eur}}$ are the US and euro risk-free rates
(Fama--French T-bill rate and one-month German bill, respectively),
and $r_{t}^{\textsc{usd/eur}}$ is the period return on the
dollar-to-euro exchange rate. All remaining factors enter the euro
panel with native European counterparts: European bank stocks,
euro-area corporate bond indices, and German government bonds in the
Duarte specification; euro-area term, credit, and inflation-linked
premia in Brooks et al.; and the 10-year Bund total-return index and
the BBB--Bund return spread in Fung and Hsieh. The Fung and Hsieh
trend-following lookback straddle factors are currency-independent by
construction (D.~Hsieh, private communication) and enter both panels
identically. The complete factor list is reported in
Appendix~\ref{app:factor_list}.

All specifications are estimated by ordinary least squares with
Newey--West HAC standard errors \citep{newey1987simple}, with the
truncation lag set by the automatic rule of \citet{newey1994automatic},
$\lfloor 4\,(T/100)^{2/9} \rfloor$, which equals four lags for every
sample in this paper. Variance inflation factors for each
specification are reported in Appendix~\ref{app:vif}. Two groups of
regressors are highly collinear---the maturity-sorted bond portfolios
$\textit{R2}$, $\textit{R5}$, $\textit{R10}$ in Duarte and the global
bond benchmarks in Brooks et al.---which inflates the standard errors
of those loadings but leaves inference on the intercept unaffected.


% ── II.B RESULTS AND DISCUSSION ───────────────────────────────────────────────

\subsection{Results and Discussion}
\label{sec:bench_results}

\paragraph{Factor Model Regressions.}

Tables~\ref{tab:duarte}, \ref{tab:fung_hsieh}, and~\ref{tab:active_fi}
report the full euro-panel regression results; the US-factor
counterparts are Tables~\ref{tab:duarte_us}, \ref{tab:fung_hsieh_us},
and~\ref{tab:active_fi_us} in Appendix~\ref{app:benchmark_us}. Across
all three frameworks, all three strategies, and both currency
denominations, the estimated intercept is positive and significant at
the $1\%$ level. The annualised alphas are economically meaningful:
$+4.0$ to $+4.6\%$ p.a.\ for the CDS--Bond basis, $+1.6$ to $+2.2\%$
for iTraxx Combined, and $+1.6$ to $+2.0\%$ for BTP~Italia.

The adjusted $R^2$ is low across all specifications, ranging from zero
to $18\%$, so the standard factor sets leave most of the return
variation unexplained---a question we take up in
Section~\ref{sec:factor_selection}. The significant loadings that do
emerge differ across strategies and are economically intuitive. The
CDS--Bond basis loads on credit and rates returns: the A-rated
industrial bond return is positive and significant in the Duarte
specification, and the 10-year government bond and credit-spread
returns are positive and significant in Fung--Hsieh, in both currency
denominations---the natural exposures of a long-credit, cash-financed
position. BTP~Italia and iTraxx Combined load instead on equity
factors, with small negative coefficients: market and size for
BTP~Italia, and size, value, and momentum for iTraxx Combined, in both
the Duarte and the Fung--Hsieh specifications. The Brooks et al.\
specification, which uses broad risk premia rather than individual
asset returns, produces fewer significant loadings but delivers
comparable alpha estimates, reinforcing that the intercepts are not an
artefact of any single factor set.

\paragraph{Cross-Model Comparison and Regime Analysis.}

Table~\ref{tab:alpha_synthesis} consolidates the alpha estimates from
the three frameworks. Panel~A reports the unconditional alphas side by
side: for each strategy, all three frameworks deliver positive and
significant alpha, and the estimates are consistent across
specifications despite their different factor structures.

The discrete regime decomposition---alpha estimated separately in the
LOW, MEDIUM, and HIGH stress regimes on the iTraxx Europe Main 5-year
spread, with the same 60 and 100 basis-point cutoffs that govern the
transaction-cost tiers of Section~\ref{sec:index_construction}
(Appendix~\ref{app:fees})---is reported for each framework in
Appendix~\ref{app:subperiod}. In every cell of those tables, alpha
rises monotonically from the LOW to the HIGH regime, and the
HIGH-regime alpha is three to seven times the LOW-regime alpha across
the nine strategy--framework combinations. The mechanism is the one
\citet{duffie2010presidential} formalises: when intermediary capital
is abundant, arbitrage forces close mispricings quickly and compress
the residual premium; when it is scarce---as it is during episodes of
widening credit spreads---arbitrageurs retreat or are forced to
unwind, mispricings widen, and the premium for those who can still
deploy capital rises. This regime dependence motivates the
cross-strategy analysis of Section~\ref{sec:interdependencies}, where
we test whether the same intermediary balance-sheet shocks drive
co-movement of mispricings across the three strategies.

\paragraph{Conditional Alpha.}

Panel~B of Table~\ref{tab:alpha_synthesis} replaces the discrete
regime split with the conditional performance evaluation framework of
\citet{christopherson1998conditioning},
in which the alpha varies continuously with a predetermined,
publicly observable state variable. We estimate the full conditional
model
\begin{equation}\label{eq:cfg}
  r_{i,t} \;=\; \alpha_0 \;+\; \alpha_1\, z_{t-1}
  \;+\; \left(\boldsymbol{\beta} + \boldsymbol{\delta}\, z_{t-1}\right)'
  \mathbf{X}_{t} \;+\; \varepsilon_{t},
\end{equation}
where $\mathbf{X}_t$ collects the factor returns of the chosen
benchmark and $z_{t-1}$ is the one-month-lagged level of the iTraxx
Europe Main 5-year spread, standardised to mean zero and unit variance
over the estimation sample. Lagging the state variable keeps it in the
investor's information set, as the conditional framework requires, and
conditioning the loadings through $\boldsymbol{\delta}$ prevents
time-varying betas from masquerading as conditional alpha. The
intercept $\alpha_0$ is then the alpha at the average stress level,
and $\alpha_1$ is the additional alpha per one-standard-deviation
increase in funding stress. Because fixed-bandwidth HAC $t$-tests
over-reject in conditional models at these sample sizes
\citep{kiefer2005new, lazarus2018har}, significance of $\alpha_1$ is
assessed by a moving-block residual bootstrap with the null imposed
($B = 9{,}999$, block length nine months); HAC inference, reported as
a robustness check, gives the same pattern.

The estimates of $\alpha_1$ in Panel~B are positive in all nine
strategy--framework cells, and the intercept $\alpha_0$ remains
significant in every case, so a baseline alpha persists at the average
stress level. For the two institutionally held strategies, $\alpha_1$
is significant in every framework: a one-standard-deviation increase
in the iTraxx Main spread adds $1.7$ to $3.1$ percentage points of
annual alpha for the CDS--Bond basis and $0.9$ to $1.5$ for iTraxx
Combined. For BTP~Italia, $\alpha_1$ is significant only under the
Duarte specification---a weaker stress sensitivity consistent with the
retail holder base that Section~\ref{sec:interdependencies} develops.
Across the nine cells, the implied alpha at $z = +1\sigma$ is two to
nine times the alpha at $z = -1\sigma$: the bulk of the unconditional
alpha is earned when funding stress is elevated, the signature of the
slow-moving capital mechanism.

\paragraph{Out-of-Sample Alpha.}

A residual concern, central to the recent literature on performance
evaluation \citep{moreira2017volatility, cederburg2020performance}, is
that the benchmark loadings in
Tables~\ref{tab:duarte}--\ref{tab:active_fi} are estimated in sample,
so the surviving alpha could reflect an in-sample-optimal combination
of the factors rather than an abnormal return an investor could have
earned in real time. Panel~C of Table~\ref{tab:alpha_synthesis}
addresses this directly. For each strategy and each benchmark, we
re-estimate the loadings on an expanding past-only window in the
recursive out-of-sample design of \citet{welch2008comprehensive}, form
the factor hedge from the lagged slopes, and define the out-of-sample
alpha as the mean of the realised hedged return
$e_{i,t} = r_{i,t} - \hat{\boldsymbol{\beta}}_{i,t-1}'\,\mathbf{f}_t$,
tested with the \citet{newey1994automatic} HAC correction. The
expanding window is initialised with a 60-month minimum, so that the
first fit has at least five observations per parameter under the
ten-factor \citet{duarte2007risk} specification, and windows with
near-singular factor designs are skipped. The out-of-sample alpha
remains positive and significant for all three strategies and all
three frameworks. The point estimates are close to the in-sample
alphas of Panel~A, with only modest attenuation, and the ranking
across strategies is preserved, with the CDS--Bond basis the
largest. The rolling-window counterpart, which additionally relaxes
the assumption of constant loadings, is reported in
Appendix~\ref{app:oos_rolling} and delivers the same conclusion.

\paragraph{Rolling Alpha.}

A natural concern with regime-based or conditional analyses is that
the estimated relationship between alpha and stress may be driven by a
single sub-period---the global financial crisis of 2008 to 2009, the
European sovereign debt crisis of 2011 to 2012, or the COVID-19 shock
of March 2020---rather than by a stable feature of the data.
Figure~\ref{fig:rolling_alpha_composite} plots rolling 36-month alpha
estimates for each strategy under each of the three benchmark
specifications. The rolling alpha is positive across essentially the
whole sample for all three strategies, so the estimates are not the
product of a single favourable sub-period, and the three benchmark
estimates move broadly together within each panel, reinforcing the
cross-model robustness of Panel~A. Detailed sub-period splits for each
framework are reported in Appendix~\ref{app:subperiod}.

Taken together, the five sets of results in this section establish two
findings that are central to the rest of the paper. First, the alpha
is robust: it survives three widely used benchmark frameworks, in both
currency denominations, in every sub-period we examine, and when each
benchmark hedge is re-estimated out of sample. Second, the alpha is
not constant: it rises systematically when intermediary capital is
scarce, whether stress enters as a discrete regime or as a continuous
predetermined state, and the high-stress alpha is a multiple of the
calm-market alpha. The first finding says the observed excess returns
are not explained by the systematic exposures the prior literature has
identified; the second points to an economic mechanism but not yet to
the specific risk factors behind the residual variation.
Section~\ref{sec:factor_selection} takes up that question on a much
larger candidate factor universe.